

\chapter{Fundamentação Teórica}\label{chap:fundamentacao}

\section{Teoria de Gêneros de Comunicação Organizacional}
\label{sec:teoria-generos}

% A fundamentação teórica deste trabalho apoia-se na tradição que trata gêneros como ação social, na formulação desenvolvida por Yates e Orlikowski ao longo de três trabalhos complementares \cite{GenresOrgComm, GenreRepertoire, GenreSystems}.

Um gênero de comunicação organizacional é definido como uma ação comunicativa tipificada, invocada em resposta a uma situação recorrente e caracterizada por \textbf{substância} e \textbf{forma} \cite{GenresOrgComm}. Onde Uma substância diz respeito aos motivos sociais, temas e propósitos expressos na comunicação e a forma suas características e a linguísticas observável. 
Os autores identificam ao menos três aspectos da forma que são relevantes: características estruturais, como listas, campos e dispositivos de estruturação da interação; o meio de comunicação empregado; e o sistema de linguagem, incluindo grau de formalidade e vocabulário especializado.

Para exemplicificar e contextualizar o que é um gênero \cite{GenresOrgComm} utilizam o exemplo de uma reunião, que é um gênero de comunicação organizacional recorrente. A situação recorrente é a demanda organizacional por interação presencial, e a ação tipificada é a reunião em si. A substância da reunião é a execução conjunta de tarefas e responsabilidades atribuídas; sua forma inclui a marcação prévia de tempo e lugar, o meio presencial e dispositivos estruturantes como a pauta e o papel de quem preside. Uma reunião específica de um comitê de pessoal em um escritório é tratada como uma \emph{instância} do gênero, e não um gênero distinto. Uma reunião não precisa ter ata nem pauta formal para ser reconhecida como reunião, mas precisa que regras distintivas suficientes identificada, dentro daquela comunidade, como instância daquele gênero \cite{GenresOrgComm}. Porém, um encontro com pessoas conversando não necessáriamente pode ser considerada uma reunião, pois não invoca regras distintivas suficientes para ser reconhecida como instância daquele gênero.

Gêneros são encenados por meio de \textbf{regras de gênero}, que associam elementos apropriados de forma e substância a determinadas situações recorrentes \cite{GenresOrgComm}. Essas regras podem operar de maneira não explícita, operando por socialização e hábito, ou padrões explícitos destinados a regular a comunicação. Um ponto decisivo para esta pesquisa é que uma instância não precisa invocar todas as regras do gênero, mas precisa invocar regras distintivas \emph{suficientes} para que a comunidade a reconheça como instância daquele gênero. O reconhecimento é uma decisão da exclusiva da comunidade, e a ele corresponde a possibilidade de reconhecimento ou não reconhecimento.

Como gêneros existem em diferentes níveis de abstração e com diferentes escopos normativos \cite{GenresOrgComm}, é possível identificar tanto tipos amplos quanto variantes especializadas dentro de uma mesma comunidade, desde que para cada uma se possa identificar uma situação recorrente, um propósito comum e características formais compartilhadas. 
% O estudo de Devitt sobre a contabilidade fiscal, discutido pelos autores, identificou seis tipos distintos de cartas e três de memorandos praticados por uma única comunidade profissional.

Comunidades raramente dependem de um único gênero. \citeauthor{GenreRepertoire} propõem a noção de \textbf{repertório de gêneros} para designar o conjunto de gêneros rotineiramente encenados por uma comunidade, cuja composição e cujo padrão de uso revelam suas práticas comunicativas e seu processo de organização. Em comunidades acadêmicas, especificamente, a produção de conhecimento é conduzida e codificada em larga medida por meio de formas de escrita, que avaliam a informação de acordo com as normas, valores e requisitos de uma disciplina. Posteriormente,em \cite{GenreSystems} apresentam o conceito de \textbf{sistemas de gêneros}: sequências de gêneros interdependentes, encenados em uma ordem típica ou em um conjunto limitado de regras aceitáveis, cujos propósitos e formas se interligam. Um sistema de gêneros estrutura seis dimensões da interação: propósito, conteúdo, participantes, forma, tempo e lugar.

\subsection{As seis dimensões da interação e a chamada de trabalhos}
\label{sec:seis-dimensoes}
 
Ao introduzir sistema de gêneros, \cite{GenreSystems} estruturam seis dimensões da interação comunicativa: propósito, conteúdo, participantes, forma, tempo e lugar . Essas dimensões funcionam como um modelo de interação sobre o qual os membros da comunidade se apoiam, e sua utilidade analítica é dupla: ao especificar o que é esperado, o sistema de gêneros especifica também, por complemento, o que não é sancionado. Ele revela quem não está autorizado a iniciar ou receber determinados gêneros, quando e onde certos gêneros não podem ser encenados, e que conteúdo ou forma é inadequado para cada um deles.
 
Aplicadas ao fluxo editorial de um simpósio, essas dimensões oferecem uma grade de análise imediata. A observação relevante para esta pesquisa é que a chamada de trabalhos, sozinha ela opera sobre praticamente todas elas, o que a torna o gênero de maior densidade estrutural do sistema. A Tabela~\ref{tab:dimensoes-cfp} sintetiza essa análise.
 
\begin{table}[H]
\centering
\caption{As seis dimensões do sistema de gêneros e sua especificação pela chamada de trabalhos}
\label{tab:dimensoes-cfp}
\footnotesize
\begin{tabular}{|p{2.6cm}|p{4.5cm}|p{6.0cm}|}
\hline
\textbf{Dimensão} & \textbf{O que estrutura} & \textbf{Como aparece na chamada} \\
\hline
Propósito (por quê) & O propósito socialmente reconhecido do sistema e de cada gênero & Pressuposto, não enunciado: a chamada indica temas, mas não explicita o que a trilha reconhece como contribuição válida \\
Conteúdo (o quê) & Quais gêneros compõem o sistema, em que sequência, e o que cada um deve conter & Lista de tópicos de interesse, escopo da trilha e elementos esperados do manuscrito \\
Participantes (quem) & Quem inicia cada gênero e a quem ele se dirige & Definição de quem pode submeter, do papel dos revisores e da coordenação, e das restrições de anonimato \\
Forma (como) & Meio, dispositivos de estruturação e sistema de linguagem & Modelo de formatação, formato de arquivo, limite de páginas e idioma \\
Tempo (quando) & Expectativas temporais associadas a cada gênero & Cadeia de prazos de submissão, notificação e versão final \\
Lugar (onde) & Localização, física ou eletrônica, de cada gênero & Sistema de submissão para a interação assíncrona e sede do evento para a apresentação \\
\hline
\end{tabular}
\end{table}


Para exemplificar o um sistema de gêneros, \cite{GenreSystems} descrevem como uma RFP(Request for Proposal) de uma organização dá origem a uma sequência de gêneros interdependentes. O fluxo de uma RFP é comumente conhecido: uma proposta é submetida por um fornecedor, este fornecedor avalia a proposta, avalia se é capaz de aceitar a proposta e elabora um relatório com uma orçamento e suas condições, desta forma informa se aceita ou se recusa participar do leilão. Normalmente este processo envolve mais de um fornecedor que por sua vez repete executa o mesmo processo e a conclusão desta negociação é a assinatura de um contrato ou a rejeição da proposta.

Desta forma é possível observar uma semelhança entre o fluxo de uma RFP e o fluxo editorial de um simpósio científico, que envolve a submissão de artigos, a avaliação por pares e a decisão final de aceitação ou rejeição. A Tabela~\ref{tab:mapeamento-generos} apresenta o mapeamento adotado, tomando como referência o Simpósio Brasileiro de Sistemas de Informação (SBSI).

A Tabela~\ref{tab:mapeamento-generos} apresenta um mapeamento baseado nestes conceitos, tomando como referência o Simpósio Brasileiro de Sistemas de Informação (SBSI).


\begin{table}[H]
\centering
\caption{Mapeamento do fluxo editorial do SBSI às categorias da Teoria de Gêneros}
\label{tab:mapeamento-generos}
\footnotesize
\begin{tabular}{|p{4.2cm}|p{9.0cm}|}
\hline
\textbf{Elemento} & \textbf{Categoria teórica} \\
\hline
Comunidade científica de SI & Comunidade que reconhece, encena e reproduz o repertório de gêneros \\
SBSI & Sistema de gêneros institucionalizado nessa comunidade \\
SBSI 2026 & Instância do sistema de gêneros, e não um gênero distinto \\
Trilha & Recorte da situação recorrente que dá origem a um subgênero do artigo, com regras próprias \\
Chamada para submissão & Gênero iniciador do sistema; fixa participantes, forma, prazo e local \\
Submissão & Gênero respondente, análogo à proposta \\
Análise preliminar & Verificação de reconhecimento de gênero pela comunidade \\
\emph{Desk rejection} & Gênero terminal de um caminho curto do sistema \\
Revisão por pares & Gênero avaliativo, análogo ao relatório de avaliação \\
Parecer & Gênero de consolidação das avaliações em recomendação \\
Aprovação ou rejeição & Gêneros terminais mutuamente exclusivos \\
\hline
\end{tabular}
\end{table}
 
A distribuição observada na Tabela~\ref{tab:dimensoes-cfp} não é homogênea e facilmente detectável, e é esta característica que interessa a este trabalho. Cinco das seis dimensões são tratadas de maneira explícita e prescritiva na chamada, em regras de gênero verificáveis por inspeção do documento submetido. A sexta, o propósito, permanece pressuposta. \citeauthor{GenreSystems} identificam o propósito socialmente reconhecido como a característica central de identificação de um sistema de gêneros, e é precisamente por ser essencial ao pertencimento à comunidade que ele tende a permanecer implícito, transmitido por enculturação e não por enunciado.
 
Uma das consequências desta análise diz respeito à natureza dos erros que levam à rejeição. Violações nas cinco dimensões codificadas são detectáveis a partir da própria chamada, pois nela estão escritas. Violações na dimensão do propósito não são, dado que um manuscrito pode estar formatado corretamente, dentro do limite de páginas, submetido no prazo e no sistema correto, e ainda assim ser rejeitado por não realizar o tipo de contribuição que aquela trilha reconhece. 

 À luz do critério de reconhecimento formulado em \cite{GenresOrgComm}, a análise preliminar não julga primariamente o mérito da contribuição: ela verifica se o \emph{gênero respondente} submetido invocou regras distintivas suficientes para ser reconhecida como pertencente ao gênero que reivindica. Escopo incompatível, desvio de formato, falta de contribuição ou problemas metodológicos conduzem à rejeição não porque o trabalho seja ruim, mas porque ele não é reconhecida como adequada a uma instância ou subgêneros. Nesses termos, a \emph{desk rejection} é a sanção da comunidade dentro de uma instância, e frequentemente decorre de uma incompatibilidade de gênero, quando o autor escreve segundo as convenções de um gênero e submete a uma trilha, a avaliação de substância, conduzida pela revisão por pares, só ocorre depois de superado esse teste de pertencimento(Triagem de diretrizes).


\subsection{Conexão com o artefato proposto}

Esse enquadramento anterior tem uma consequência direta para o propósito do artefato desenvolvido nesta pesquisa. Dentro do que propôe a Teoria de Gêneros de Comunicação Organizacional, a dificuldade do pesquisador iniciante pode não ser necessariamente de escrita nem de método, existe a possibilidade de ainda não haver internalizado o repertório de gêneros de uma comunidade. Tal internalização ocorre por enculturação, processo tipicamente lento, tácito e dependente de exposição prolongada às práticas do grupo \cite{GenreRepertoire}. Como as regras de gênero de cada trilha não tendem a ser enunciadas de forma completa, o autor iniciante dispõe apenas de indícios parciais, distribuídos entre a chamada de trabalhos, os artigos já publicados e os pareceres que eventualmente recebe.

%CITAR ALGO RELACIONADO A DESK REJECTION.

O artefato proposto neste trabalho é conceitualmente um instrumento que acelera e tutora o autor iniciante nesta enculturação, tornando operáveis regras que de outro modo permaneceriam subentendidas. Trata-se, portanto, de uma contribuição de natureza teórica e não apenas tecnológica. Assim decorrem as três questões que orientam o desenho da solução: quais são as regras de gênero de cada trilha, como um autor pode verificar se seu manuscrito está em conformidade com elas, e de que modo um sistema multiagente pode realizar essa verificação de maneira fundamentada.

\section{Análise de \textbf{moves} retóricos}

A Teoria de Gêneros fornece o porquê sociotécnico, mas é abstrata demais para orientar diretamente a implementação. A ponte entre o construto abstrato e o texto observável é a análise de \textbf{moves} retóricos, na tradição da escrita acadêmica de \cite{GenreAnalysis}. O modelo CARS (\textbf{Create a Research Space}) decompõe a introdução de artigos científicos em movimentos retóricos recorrentes: estabelecer o território (situar a relevância e revisar a literatura), estabelecer o nicho (apontar a lacuna ou o problema) e ocupar o nicho (anunciar a contribuição e a organização do trabalho).

A importância desse modelo para a pesquisa é dupla. Primeiro, os \textbf{moves} são a manifestação observável e mensurável do gênero abstrato: onde \cite{GenresOrgComm} falam de regularidades de forma e substância, \cite{GenreAnalysis} oferece o microscópio que as torna visíveis no texto. Segundo, os \textbf{moves} são operacionalizáveis por processamento de linguagem natural, o que fornece precedente metodológico para sua extração automatizada. A articulação, portanto, é a seguinte: a Teoria de Gêneros define o construto de Sistemas de Informação; a análise de \textbf{moves} operacionaliza esse construto em unidades observáveis; 
e a extração automatizada de \textbf{moves} fornece o mecanismo pelo qual um modelo de LLM pode verificar se uma instância invocou regras distintivas suficientes para ser reconhecida como pertencente ao gênero que reivindica.

\section{Design Science Research e a anatomia da teoria de design}
%(HEVNER et al., 2004; PEFFERS et al., 2007)
O trabalho adota o método Design Science Research \cite{Hevner2004,Peffers2007}, que orienta a construção e a avaliação de artefatos como forma de produção de conhecimento. Mais do que produzir um artefato, esta pesquisa pretende produzir uma \textbf{teoria de design}, no sentido de \cite{DesignTheory}, cuja anatomia prevê oito componentes: propósito e escopo; construtos; princípios de forma e função; mutabilidade do artefato; proposições testáveis; conhecimento justificador; princípios de implementação; e instanciação expositiva. 
%Essa anatomia estrutura o Capítulo 4 e é o que permite que a contribuição transcenda o protótipo específico do SBSI, elevando-a a um conjunto de princípios generalizáveis — condição para a evolução do trabalho rumo a uma tese de doutorado.

% --- Preâmbulo (já incluídos para a Figura 1) ---
% \usepackage{tikz}
% \usetikzlibrary{positioning, fit, backgrounds, arrows.meta}

\begin{figure}[H]
  \centering
  \begin{tikzpicture}[
    >={Stealth[round]},
    font=\small,
    teoria/.style={draw, rounded corners, align=left,
                   text width=6.2cm, inner sep=8pt, minimum height=1.7cm},
    lateral/.style={draw, rounded corners, align=left, fill=gray!4,
                    text width=3.0cm, inner sep=8pt, minimum height=1.7cm},
    nivel/.style={font=\footnotesize},
    setinha/.style={->, thick, gray}
  ]

    % --- Camada 1: fenômeno ---
    \node[teoria, fill=violet!8] (t1)
      {\textbf{Teoria de Gêneros}\\[2pt]{\footnotesize Yates e Orlikowski}\\{\footnotesize \emph{fenômeno de SI}}};
    \node[lateral, right=0.7cm of t1] (l1)
      {\textbf{o quê?}\\[2pt]{\footnotesize $\rightarrow$ base de conhecimento}};

    % --- Camada 2: operacionalização ---
    \node[teoria, fill=teal!10, below=0.8cm of t1] (t2)
      {\textbf{Análise de moves --- Swales}\\[2pt]{\footnotesize modelo CARS}\\{\footnotesize \emph{operacionalização}}};
    \node[lateral, right=0.7cm of t2] (l2)
      {\textbf{como medir?}\\[2pt]{\footnotesize $\rightarrow$ extração por artigo}};

    % --- Camada 3: arquitetura (kernel) ---
    \node[teoria, fill=orange!12, below=0.8cm of t2] (t3)
      {\textbf{Teoria de compiladores}\\[2pt]{\footnotesize kernel computacional}\\{\footnotesize \emph{justificação da arquitetura}}};
    \node[lateral, right=0.7cm of t3] (l3)
      {\textbf{por que funciona?}\\[2pt]{\footnotesize $\rightarrow$ verificação e tutor}};

    % --- Setas entre as camadas ---
    \draw[setinha] (t1) -- (t2);
    \draw[setinha] (t2) -- (t3);

    % --- Moldura do DSR englobando tudo ---
    \begin{scope}[on background layer]
      \node[draw, dashed, rounded corners, fill=gray!2,
            fit=(t1)(l1)(t3)(l3), inner sep=18pt] (dsr) {};
    \end{scope}
    \node[anchor=north west, font=\bfseries]
      at ([shift={(4pt,-4pt)}]dsr.north west) {Método};

  \end{tikzpicture}
  \caption{Integração das três camadas teóricas sob o método DSR. Elaboração própria.}
  \label{fig:integracao-teorias}
\end{figure}

% > **[Figura 2 — Quadro integrador das três teorias]**
% > *Melhor local para inserção: ao final do Capítulo 2, como síntese. Diagrama em três camadas (Teoria de Gêneros → análise de moves → teoria de compiladores) dentro da moldura do DSR, cada camada associada à pergunta que responde (o quê / como medir / por que funciona) e ao estágio correspondente do artefato.*
% > 
% > Figura 2. Integração das três camadas teóricas sob o método DSR. Elaboração própria.



% ------------------------------------------------------------------------------------------------------------

% \section{Contextualization}

% Entanglement Theories seek to understand the phenomena that emerge from interactions between humans and non-humans \cite{frauenberger2019entanglement}. LAKS JJKSF JSJKSD ssk sdjfksd fsdjk dsjdks sdjkf jksjks fdjksdfj sjsjsdjfksd fsdjk dsjdks sdjkf jksjks fdjksdfj sjsjsdjfksd fsdjk dsjdks sdjkf jksjks fdjksdfj sjsjsdjfksd fsdjk dsjdks sdjkf jksjks fdjksdfj sjsjsdjfksd fsdjk dsjdks sdjkf jksjks fdjksdfj sjsjsdjfksd fsdjk dsjdks sdjkf jksjks fdjksdfj sjsjsdjfksd fsdjk dsjdks sdjkf jksjks fdjksdfj sjsj. Table \ref{tab:fontes_incerteza_tar} presents bla blab alb al.

% \begin{table}[ht!]
% \centering
% \scriptsize
% \caption{Sources of Uncertainty in ANT.}
% \rowcolors{3}{gray!10}{white}
% \begin{tabular}{p{0.19\textwidth} p{0.32\textwidth} p{0.40\textwidth}}
% \toprule
% \textbf{Soursdfsinty} & \textbf{Descsdfsdtion} & \textbf{Impliss} \\
% \midrule

% No Groxxdfsup vesdfsdfsdation: & Groups asdfdfsfsdfsdvisional assemblages. & The researcher must empiricasdfsdfsdfsdsd transformed within networks. \\

% Isolated sdfsdfsd  fsd f & sdfsd fsdf sdf sd fsddf sd sdfs fsdf dssd fs fsfsdfd sds dfsd fs. & sf sd fdf ds fds fsd sdfdsf sdf sd fsd fsd sdf sd. \\

% Human s jfsf sdfsd human Agency & sdj fskjf jkfskjfsdjksdfks jkfdsjks fjks sfkdfssfdkksskdkssdk. & sjf sdkjfsdf kjskjfsjffsfsfsdfsdfsdjkfj fkf kjs fdkfskfskfskjskjfsfsk. \\

% sd fljkfsdjk fsdjkfsdjkfsdjkfsdkjsfdkjs & s dfjksfjksfjk sdfjksd jksjksjkfskd sdk sksksksksk. & sdjkf sdjkfsfkjsdfkj sdjkfkssfkfsdkjfsdkksfksds. \\

% \bottomrule
% \end{tabular}
% \label{tab:fontes_incerteza_tar}
% \caption*{\textbf{Source:} Author's elaboration.}
% \end{table}




% The first source of uncertainty \cite[pp.~27-42]{latour2005reassembling} statkba jkdajb dakb jadjkbkadjb ajkbdajkmation, the researcher cannot simply rely on fixed categories such as ``society,'' ``organization,'' or ``users'' as starting points. It becomes necessary to observe how sfk jsfkjsd fjks fjkssfgroupings.

% The second source of uncertainty \cite[pp.~44-62]{latour2005reassembling} concerns the nature of action: action is always distributed. There is no isolated agent or clear point of origin, since ev
% This way of thinking about action profoundly transforms how we investigate social and technical phenomena. If action is always distributed, the researcher can no longer attribute it to a single agent such as a person, an institution, or a technology. Instead, it becomes necessary to trace the trajectory of the action and observe how it is produced by a network of heterogeneous elements that operate together. This requires close attention to detail: which tools were used, which norms were involved, which interpretations were mobilized, and which artifacts intervened. Agency ceases to be something that someone possesses and becomes something that occurs among many.
